% Generated draft: requires mathematical and source review.
\documentclass[11pt]{article}
\usepackage[margin=25mm]{geometry}
\usepackage{fontspec}
\setmainfont{Latin Modern Roman}
\usepackage{amsmath,amssymb,mathtools,unicode-math}
\setmathfont{Latin Modern Math}
\usepackage{enumitem,needspace,xurl,hyperref}
\setlength{\parindent}{0pt}
\setlength{\parskip}{5pt}
\setlength{\emergencystretch}{3em}
\newcommand{\N}{\mathbb{N}}
\newcommand{\Z}{\mathbb{Z}}
\newcommand{\Q}{\mathbb{Q}}
\newcommand{\R}{\mathbb{R}}
\newcommand{\C}{\mathbb{C}}
\newcommand{\F}{\mathbb{F}}
\newcommand{\A}{\mathbb{A}}
\newcommand{\PP}{\mathbb{P}}
\newcommand{\E}{\mathbb{E}}
\newcommand{\eps}{\varepsilon}
\newcommand{\dd}{\,\mathrm{d}}
\begin{document}
\section*{UnsolvedMath ID: 30001463}
\textbf{Source problem code:} catalogue entry
\subsection{Definitions and mathematical statement}
Let \(t\in\{\widetilde{A}_2,\widetilde{C}_2\}\). Write \(\mathcal D_t\) provisionally for the collection of combinatorial data admitted by the construction referred to in the supplied record, and write \((X_d,\Gamma_d)\) for the building and panel-regular lattice produced from \(d\in\mathcal D_t\). These are placeholders: the record does not specify the admissible data, the construction, or its auxiliary choices, so it does not determine \(\mathcal D_t\) or the assignment \(d\mapsto(X_d,\Gamma_d)\).

A panel in a two-dimensional simplicial building is a codimension-one face of a chamber, hence an edge. A regular action on a specified set means that for any two elements of that set there is exactly one group element taking the first to the second. The record does not specify the panel types on which regularity is required, whether the action preserves types, or the precise automorphism group in which each lattice is taken.

\textbf{Source problem, conditional on specifying the construction.} For the constructions of each indicated type, determine whether distinct admissible data can yield isomorphic buildings, and determine when this occurs. In provisional notation, the existence question is
\[
\exists\,d,e\in\mathcal D_t\quad
\bigl(d\ne e\ \text{ and }\ X_d\cong X_e\bigr).
\]
Likewise, determine whether distinct admissible data can yield commensurable lattices, and determine when this occurs. Finally, determine which of the buildings \(X_d\) are classical and which are exotic. These are separate questions; the supplied record does not assert that building isomorphism and lattice commensurability are equivalent.

The terms in these questions require further specification. Abstract commensurability of groups \(\Gamma\) and \(\Lambda\) means
\[
\exists\,H\leq\Gamma\ \exists\,K\leq\Lambda:
[\Gamma:H]<\infty,\quad[\Lambda:K]<\infty,
\quad H\cong K.
\]
An ambient interpretation instead requires specified embeddings into a common group \(G\); commensurability up to conjugacy then asks for some \(g\in G\) such that \(\Gamma\cap g\Lambda g^{-1}\) has finite index in both \(\Gamma\) and \(g\Lambda g^{-1}\). The record does not select between these interpretations or supply such embeddings. It also does not state whether building isomorphisms must preserve vertex types, or whether distinct data means literal inequality or inequivalence under specified changes of presentation. Under a usual terminology, a classical affine building is a Bruhat--Tits building and an exotic building is one outside that class; the intended class and hypotheses must be checked against the construction. Consequently, the displayed formula is a conditional formulation of the source question, not a fully specified existence assertion about a supplied collection of data.
\subsection{Short English statement}
For the cited constructions of panel-regular lattices in buildings of types \(\widetilde{A}_2\) and \(\widetilde{C}_2\), determine when different combinatorial inputs yield isomorphic buildings or commensurable lattices, and which resulting buildings are classical or exotic. The record omits the construction and several equivalence conventions, so a complete precise formulation requires further source information. The supplied open status is a dated curation claim.
\subsection{Sources}
\begin{itemize}
\item Local UnsolvedMath registry, \texttt{data/unsolvedmath/problems.json}, permanent ID \(30001463\), source problem code \texttt{catalogue entry}, entitled \emph{Isomorphism and Commensurability of Panel-Regular Building Lattices}; supplied complete source record in the Oberwolfach Reports Open Problems collection.
\item \emph{Non-positive Curvature and Geometric Structures in Group Theory} (2010), Oberwolfach Reports, pp.\ 1165--1224, DOI \texttt{10.4171/owr/2010/20}. Supplied report reference and literature source labelled ``original OWR report'': \url{https://doi.org/10.4171/owr/2010/20}. The background describes these reports as documenting mathematical workshops and their open-problem discussions. No individual proposer or precise problem location within the report is supplied; the proposed year is recorded as 2010.
\item Supplied literature assessment, checked \texttt{2026-08-21}: status ``open.'' The assessment says that the source explicitly leaves deciding when two data sets yield isomorphic buildings or commensurable lattices open, and that no general resolution was found. This is supplied dated curation, not an independently verified literature assessment. The supplied statement also labels the classical-versus-exotic question open. The background expressly requests independent verification of time-sensitive status and cited claims.
\item The registry labels the statement \texttt{corrected\_verified} and records a manual semantic, typographic, and LaTeX review asserting that its cleanup is self-contained and adds no unsupported source-dependent content. These are supplied review labels; they do not resolve the missing construction details or constitute mathematical verification of this draft.
\end{itemize}
\end{document}
